\section{Notation and diffusion conventions}
\label{app:notation}

\subsection{Symbols}

\begin{center}
\small
\begin{tabular}{ll@{\qquad}ll}
\toprule
symbol & meaning & symbol & meaning \\
\midrule
$\nsites$ & chain length $=$ ambient dimension & $\kernel,\ \kernelth$ & transition kernel \\
$\nseq$ & number of observed sequences & $\innov$ & innovation density \\
$\ngrid$ & number of quadrature grid points & $\corr$ & autoregressive coefficient \\
$\gridpt{1},\dots,\gridpt{\ngrid}$ & grid points & $\ivar$ & innovation variance \\
$\gridh$ & grid spacing & $\ncomp$ & mixture components \\
$\halfwidth$ & grid half-width, domain $[-\halfwidth,\halfwidth]$ & $\mu$ & initial law $\mu(a_1)$ \\
$t$ & diffusion time & $\like$ & unary likelihood factor \\
$\Dt = 1-e^{-2t}$ & noise variance & $\lmsg,\ \rmsg$ & forward, backward message \\
$\Pz,\ \Pt$ & clean, noised distribution & $\params$ & kernel parameters \\
$\score(x,t)$ & score $\nabla_x\log\Pt(x)$ & $\Qfun$ & EM auxiliary function \\
$\pmean{a}$ & posterior mean of $a$ given $x$ & $\Xis$ & expected transition statistic \\
$\gwn$ & standardised Gaussian white noise & $\Lik$ & marginal log-likelihood \\
$\tmin,\ \tmax$ & reverse-integration endpoints & $\lmmse$ & LMMSE estimator \\
\bottomrule
\end{tabular}
\end{center}

Two symbols carry a history worth stating, because an earlier draft of this work used them
differently. \textbf{$\kernel$ always means the transition kernel} and never a grid size; the
grid size is always $\ngrid$. The ambient dimension is written $\nsites$ because it \emph{is} the
chain length and every graphical-model statement is indexed by sites; the diffusion literature
more often writes $n$ for the sample count, which here is $\nseq$.

\subsection{Time orientation and the reverse SDE}

\begin{center}
\fbox{\parbox{0.93\textwidth}{
\textbf{Convention, used without variation.} $\gwn$ is standardised Gaussian white noise,
$\langle\gwn(t)\gwn(t')\rangle=\delta(t-t')$, so every $\sqrt{2}$ is written explicitly.

\emph{Forward}, $t$ running from $0$ to $\tmax$:
\[
\frac{\mathrm{d}x}{\mathrm{d}t} = -x + \sqrt{2}\,\gwn(t),
\qquad
x(t) = e^{-t}a + \sqrt{\Dt}\,z,
\qquad
\Dt = 1-e^{-2t},\quad z\sim\normal{0}{I}.
\]

\emph{Reverse}, in the elapsed-time variable $\tau = \tmax - t$ running from $0$ to
$\tmax-\tmin$, so that increasing $\tau$ means decreasing $t$:
\[
\frac{\mathrm{d}x}{\mathrm{d}\tau} = x + 2\,\score\!\left(x, \tmax-\tau\right) + \sqrt{2}\,\gwn(\tau).
\]

\emph{Probability-flow ODE}, same marginals, deterministic:
$\;\mathrm{d}x/\mathrm{d}\tau = x + \score(x,\tmax-\tau)$.

Initialisation $x(\tau{=}0)\sim\normal{0}{I}$; termination at $\tau = \tmax-\tmin$.
}}
\end{center}

The reverse drift follows from the time-reversal theorem for diffusions
\citep{anderson1982,haussmann1986}: for a forward SDE $\mathrm{d}x = f\,\mathrm{d}t +
g\,\mathrm{d}W$ the time-reversed process has drift $f - g^2\nabla\log p_t$ and the same
diffusion coefficient. With $f=-x$ and $g=\sqrt{2}$ the backward-in-$t$ drift is
$-x - 2\score$, and rewriting in $\tau$ flips the sign to give the box above. Because unit-variance
data keep $\Var(x_t)=1$, initialising at $\normal{0}{I}$ incurs a bias $e^{-2\tmax}$, which is
$2.5\times10^{-3}$ at $\tmax=3$. The implementation is
\codefile{src/reverse.py}, whose module docstring states the same convention; the sampler uses a
geometric time grid, denser at small $t$, for the stiffness reason given in
\S\ref{sec:generation}.

\section{The score--posterior identity}
\label{app:identity}

\paragraph{Differentiation under the integral.}
$\Pt(x) = \int\mathrm{d}a\,\Pz(a)\,g_{\Dt}(x-e^{-t}a)$ with
$g_{\Dt}(u) = (2\pi\Dt)^{-\nsites/2}e^{-\|u\|^2/2\Dt}$. Since
$\nabla_x g_{\Dt}(x-e^{-t}a) = -\frac{x-e^{-t}a}{\Dt}g_{\Dt}(x-e^{-t}a)$ and, for $x$ in a
compact neighbourhood, $\|x-e^{-t}a\|\,g_{\Dt}(x-e^{-t}a) \le c_1 e^{-c_2\|a\|^2}$ uniformly, the
integrand is dominated by an integrable function whenever $\Pz$ has a finite first moment.
Dominated convergence then licenses
\[
\nabla_x \Pt(x)
= \int \mathrm{d}a\; \Pz(a)\,\frac{e^{-t}a - x}{\Dt}\, g_{\Dt}(x-e^{-t}a).
\]
Dividing by $\Pt(x)$ and recognising $\Pz(a)g_{\Dt}(x-e^{-t}a)/\Pt(x) = P(a\mid x,t)$ gives
\eqref{eq:score}. Nothing is approximated. In the empirical-Bayes literature this is Tweedie's
formula \citep{robbins1956,miyasawa1961,efron2011}.

\paragraph{Risk decomposition.} For any $u$ that is a function of $x$ alone,
\begin{equation}
\label{eq:orthogonal}
\E\!\left[\|a-u\|^2 \,\middle|\, x\right]
= \E\!\left[\|a - \pmean{a}\|^2 \,\middle|\, x\right] + \|u - \pmean{a}\|^2 ,
\end{equation}
by expanding $a-u = (a-\pmean{a}) + (\pmean{a}-u)$ and noting the cross term vanishes because
$\pmean{a}-u$ is $x$-measurable and $\E[a-\pmean{a}\mid x]=0$. Hence $\pmean{a}$ minimises the
conditional risk, with an irreducible floor given by the first term \citep{lehmann1998}.
Optimality is with respect to squared error specifically; a different loss selects a different
functional of the posterior. This is why every table in \S\ref{sec:pointwise} reports error
\emph{against} $\pmean{a}$ rather than raw risk: the floor is common to all estimators and
subtracting it is what makes the comparison informative.

\section{Sum-product on the posterior chain}
\label{app:bp}

\paragraph{Posterior factorisation.} Bayes' rule applied to \eqref{eq:prior}, with the likelihood
factorising sitewise because the forward noise is independent across coordinates, gives
\eqref{eq:posterior}. Each likelihood factor has degree one among latent variables, so attaching
it to the prior's path adds a leaf and cannot create a cycle: the posterior factor graph is a
tree. Conditioning usually creates dependence; it does not here, precisely because each
observation is generated from a single latent variable. The construction fails the moment the
forward noise is correlated across coordinates, since then $P(x\mid a)$ does not factorise, the
likelihood becomes a factor touching several $a_i$, and the posterior graph acquires edges.

\paragraph{Exactness on a tree.} Let the graph be a tree and let $a_i$ be a variable node.
Removing $a_i$ disconnects it into components, and every factor involves $a_i$ together with
variables of exactly one component --- a factor touching two components would create a second
path between them, hence a cycle. Grouping factors by component gives
$P \propto \prod_c g_c(a_i, v^{(c)})$, so conditionally on $a_i$ the components are independent.
Unrolling the forward recursion \eqref{eq:messages},
\[
\lmsg_i(a_i) \propto \int\!\mathrm{d}a_{1:i-1}\;
\mu(a_1)\prod_{j=2}^{i}\kernel(a_j\mid a_{j-1})\prod_{j=1}^{i-1}\like(x_j\mid a_j),
\]
which is the joint density of everything at or left of site $i$ with the left variables
integrated out and $\like(x_i\mid a_i)$ \emph{not} included; symmetrically for $\rmsg_i$. The two
are conditionally independent given $a_i$, so their product times the local evidence is
proportional to the marginal. The pairwise statement is the same argument with the separator
taken to be the edge.

\paragraph{Where the local factor goes.} In \eqref{eq:messages} the local evidence at site $i-1$
is absorbed \emph{before} the transition, so $\lmsg_i$ carries the factors
$\like(x_1\mid a_1),\dots,\like(x_{i-1}\mid a_{i-1})$ and $\rmsg_i$ carries
$\like(x_{i+1}\mid a_{i+1}),\dots,\like(x_{\nsites}\mid a_{\nsites})$, and each appears exactly
once in the marginal. Absorbing it \emph{after} would place $\like(x_i\mid a_i)$ in both
messages and square it: the
belief would be sharper than the truth and every posterior variance too small. This produces
plausible-looking output, which is why the convention is pinned by
\codefile{tests/test\_message\_normalization.py} rather than left to a comment.

\paragraph{Evidence.} Messages are renormalised at each step to prevent underflow; accumulating
the discarded constants recovers $\log\Pt(x)$ for the discretised model exactly. This is the
objective that \S\ref{sec:em} maximises, and having it in closed form is what makes it possible
to check a gradient against a finite difference of the thing being optimised.
\coderef{src/bp\_grid.py}{grid\_bp} returns it alongside the beliefs;
\coderef{src/exact\_scores.py}{gaussian\_log\_evidence} is the closed form it is tested against.

\paragraph{Complexity at the functional level.} Each message update integrates one variable
against the kernel. Abstractly this is an integral operator applied to a function; represented on
$\ngrid$ points it is one $\ngrid\times\ngrid$ matrix--vector product, $O(\ngrid^2)$. There are
$2\nsites$ updates, so one sequence costs $O(\nsites\ngrid^2)$ against $O(\ngrid^{\nsites})$ for
naive marginalisation. At $\nsites=32,\ngrid=401$ that is about $5\times10^6$ operations.
Batching over sequences turns each update into a single GEMM, which is why the same code runs on
a GPU without restructuring (\codefile{src/backend.py}, gated by
\codefile{tests/test\_backend\_parity.py}).

\section{The Gaussian model and the second-order baseline}
\label{app:lmmse}

For $\innov=\normal{0}{\ivar}$ with $\ivar=1-\corr^2$ and $a_1\sim\normal{0}{1}$, the vector $a$
is jointly Gaussian with $(\Sigma_0)_{ij}=\corr^{|i-j|}$. Expanding
$\Pz(a)\propto\exp[-\tfrac{1}{2\ivar}\sum_{i\ge2}(a_i-\corr a_{i-1})^2 - \tfrac12 a_1^2]$
produces only $a_i^2$ and $a_ia_{i-1}$ terms, so $\Sigma_0^{-1}$ is tridiagonal with
$(\Sigma_0^{-1})_{ii} = (1+\corr^2)/\ivar$ in the interior, $1/\ivar$ at the two ends, and
$-\corr/\ivar$ on the off-diagonals. Adding the likelihood contributes $(e^{-2t}/\Dt)I$, so the
posterior precision is tridiagonal at every $t$: Markovianity is a statement about the precision,
while the covariance $\Sigma_0$ is dense.

Since $x = e^{-t}a + \sqrt{\Dt}z$ with $z$ independent, $(a,x)$ is jointly Gaussian with
$\Cov(a,x) = e^{-t}\Sigma_0$ and $\Cov(x) = e^{-2t}\Sigma_0 + \Dt I =: \Sigma_t$, whence
$\E[a\mid x] = e^{-t}\Sigma_0\Sigma_t^{-1}x$ and, by \eqref{eq:score},
$\score(x,t) = -\Sigma_t^{-1}x$. For zero-mean jointly Gaussian variables the conditional
expectation is linear, so it coincides with the LMMSE estimator \citep{kay1993}, which proves the
second half of Proposition~\ref{prop:lmmse}. The first half is the observation that both maps in
\eqref{eq:closure} read only the mean and variance of $\innov$, so the whole recursion is a
function of $(0,\ivar)$ alone and therefore returns what it returns on the Gaussian chain.

Consequently, on this model, the \emph{message-representation error} and the \emph{model error}
coincide: projecting messages onto Gaussians and replacing the prior by its covariance-matched
Gaussian give the same estimator. This is a property of single-Gaussian closure with linear
transitions, not of Gaussian approximations in general, and the two separate for richer message
families.

\section{Numerical representation}
\label{app:numerics}

Messages are represented on a uniform grid $\gridpt{1},\dots,\gridpt{\ngrid}$ spanning
$[-\halfwidth,\halfwidth]$ with spacing $\gridh = 2\halfwidth/(\ngrid-1)$ and composite
trapezoidal weights $w_k = \gridh$ for interior points, $\gridh/2$ at the two ends. All
recursions run in the log domain with a per-step maximum subtracted before exponentiating, and
the removed constants are restored in the evidence, so no message underflows.

\paragraph{Two error sources.} \emph{Truncation}: probability mass outside
$[-\halfwidth,\halfwidth]$ is discarded rather than transported, an error controlled by
$\halfwidth$ and essentially independent of $\ngrid$. \emph{Quadrature}: the trapezoidal rule is
second order in $\gridh$ for integrands with bounded second derivative, an error controlled by
$\gridh$ at fixed $\halfwidth$. Figure~\ref{fig:grid} measures them separately, using the
worst-case boundary mass over sites for the first and the interior column-mass residual of the
transition matrix for the second; the interior restriction $|\gridpt{}| \le \halfwidth/2$ removes
columns whose transition density has a tail outside the grid, which is what makes the second
diagnostic measure quadrature alone. The observed rates are: boundary mass
$3\times10^{-4} \to 1.4\times10^{-8} \to 1.4\times10^{-10}$ at $\halfwidth = 4, 8, 10$; interior
residual falling by a factor of four per halving of $\gridh$, i.e.\ $O(\gridh^2)$.

At $\halfwidth=4$ the interior residual saturates at $\sim10^{-3}$ instead of continuing to fall,
because at that domain even the ``interior'' columns have truncated tails --- the floor is the
truncation error leaking into the quadrature diagnostic, and it is visible in the data
(\codefile{outputs/exp\_18/boundary.csv}). This is the reason the two diagnostics have to be read
together and neither alone certifies convergence.

\paragraph{What is not claimed.} A grid-resolution sweep at fixed $\halfwidth$ does not rule out
truncation error, and a domain sweep at fixed $\ngrid$ does not rule out quadrature error. We
report both axes and the two diagnostics; we do not claim continuum convergence, only that at
$\ngrid=401$, $\halfwidth=8$ both diagnostics sit far below every effect the paper reports.
For the Laplace family, whose innovation density has a discontinuous derivative at the origin, we
claim second order and no better.

\section{Learning procedure}
\label{app:learning}

\paragraph{Objective.} Maximum marginal likelihood
$\Lik(\params) = \sum_\mu\log P_{t,\params}(x^\mu)$, with the complete-data objective
$\log P_\params(a,x) = \log\mu(a_1) + \sum_{i\ge2}\log\kernelth(a_i\mid a_{i-1}) +
\sum_i\log\like(x_i\mid a_i)$. Only the middle term depends on $\params$, which is why $\Xis$ of
\eqref{eq:xi} is sufficient and why $\Qfun(\params\mid\params') = \langle\Xis(\params'),
\log\kernelth\rangle$ up to $\params$-free constants.

\paragraph{E-step.} One forward and one backward sweep per sequence, then a rank-one accumulation
of the pairwise statistics: writing $\ell^{(i)}$ for the vector $\like(x_i\mid \cdot)$ on the
grid, and with $f_i = w\odot\lmsg_i\odot\ell^{(i)}$ and $g_i = w\odot\rmsg_i\odot\ell^{(i)}$,
the pairwise belief on edge $(i,i+1)$ is
$f_i(j)\,\kernel[k,j]\,g_{i+1}(k)/Z_i$ with $Z_i = g_{i+1}^\top \kernel f_i$, and
$\Xis = \left(\sum_{\mu,i} g_{i+1}(f_i)^\top/Z_i\right)\odot \kernel$. $\Xis$ sums to the number
of edges in the dataset, which is asserted in the tests. Exact for the discretised model, in the
sense of Remark~\ref{rem:exact}(ii). Implementation \coderef{src/em.py}{e\_step}.

\paragraph{M-step, by kernel.}
\begin{itemize}\itemsep2pt
\item \emph{Gaussian}, $\params=(\corr,\ivar)$: closed-form maximiser, the belief-weighted
      Yule--Walker equations $\corr = S_{01}/S_{00}$, $\ivar = (S_{11}-\corr S_{01})/W$. Exact
      maximisation, hence exact EM for that family.
\item \emph{Laplace}, $\params=(\corr,b)$: closed form despite non-differentiability. Profiling
      $b$ out of $\Qfun$ turns the problem into minimising the belief-weighted mean absolute
      residual, whose exact solution is a weighted median. Exact maximisation.
\item \emph{Mixture} \eqref{eq:kernel}: the complete-data problem carries a second latent
      variable, so the M-step is itself an inner EM over the mixture label with
      $\Xis$ held fixed. We run \textbf{four} inner conditional-maximisation sweeps, alternating
      the $(\pi,\nu,s^2)$ block and the $\corr$ block, both closed form. This increases $\Qfun$
      without maximising it: \textbf{generalised EM / ECM} \citep{meng1993}, not exact EM.
\item \emph{Mixture-density network}: one backtracked Adam step, accepted only if $\Qfun$
      increased, halving the step up to five times otherwise and standing still if no ascent is
      found. Also generalised EM.
\end{itemize}

\paragraph{Guarantees, initialisation, stopping.} Generalised EM guarantees monotone ascent of
$\Lik$ \citep{dempster1977,wu1983}; convergence is to a stationary point, not necessarily a global
maximum. Mixture components are initialised spread over $[-\sqrt{\Var}, \sqrt{\Var}]$ with random
scales; $\corr$ is initialised away from its truth in every experiment --- at $0.3$ in most, and
over $\{0,\,0.3,\,0.6,\,-0.4\}$ in the initialisation sweep of \codefile{exp\_18}. The truth is
$0.85$ in \codefile{exp\_02}, \codefile{exp\_16}, \codefile{exp\_18} and \codefile{exp\_22}, and
$0.8$ in \codefile{exp\_06}, \codefile{exp\_07} and \codefile{exp\_08}; every table and figure
states the value it used. Iteration
stops when the relative increase of $\Lik$ falls below $10^{-9}$, capped at $50$--$120$
iterations depending on the experiment. Variances are floored at $10^{-6}$ to prevent component
collapse; no other regularisation is applied, and the component means are unconstrained
(\S\ref{sec:em}). Monotonicity is recorded per run as the largest decrease across iterations and
is zero in every run reported here.

\section{Experimental protocol}
\label{app:protocol}

\begin{center}
\small
\begin{tabular}{ll}
\toprule
chain & Laplace AR(1), $\corr=0.85$, $\ivar=1-\corr^2$, $a_1\sim\normal{0}{1}$, $\nsites=32$ \\
other families & Gaussian, Student-$t$ ($\nu=5$), uniform, symmetric two-component mixture
                 ($\kappa=0.6$) \\
                & all variance-matched to $\ivar$, so $\Cov(a_i,a_j)=\corr^{|i-j|}$ exactly \\
grid & $\ngrid=401$, $\halfwidth=8$ ($\gridh=0.04$); refinements $\ngrid\in\{801,1601\}$ \\
noise levels & $t\in\{0.1,0.2,0.4,0.8,1.6\}$ for training and pointwise evaluation \\
budgets & $\nseq\in\{32,128,512,1024,2048,4096\}$ \\
test set & $256$ held-out sequences, seed \texttt{exp07-test}, disjoint from all training seeds \\
replicates & six; the replicate index is mixed into every training seed, so replicates vary \\
           & the training set \emph{and} the initialisation; the test set is common to all six \\
\midrule
EM--BP & mixture-innovation kernel \eqref{eq:kernel}, $\ncomp=4$ unless swept;
         $12$ free parameters \\
       & $\corr$ initialised at $0.3$; $50$--$120$ EM iterations; tolerance $10^{-9}$ \\
global MLP & fully connected, $25{,}248$ parameters, denoising score matching, $20{,}000$
             steps \\
           & both parameterisations ($\varepsilon$ and $a$) trained; Adam \\
local CNN & weight-shared one-dimensional convolution, radius $r$, receptive field $2r+1$ \\
capacity control & $24$ configurations, $3.2$k--$297$k parameters \\
\midrule
hardware & CPU: single-threaded, \texttt{medium\_cpu} partition, Bocconi HPC \\
         & GPU: CUDA 12.4, CuPy 13.6, gated by a CPU/GPU parity test before any sweep \\
\bottomrule
\end{tabular}
\end{center}

\paragraph{Selection rules, stated plainly.} There is \textbf{no validation split}. Two selections
therefore happen on the evaluation set and both favour the baselines: (i) Table~\ref{tab:pointwise}
reports, at each noise level, the better of the two denoising-score-matching parameterisations,
which is oracle post-selection over two models; (ii) the convolution's receptive field is a fixed
configuration value read off the same sweep that is evaluated on the test set. The numbers in
\S\ref{sec:pointwise} that depend on (ii) are labelled exploratory. Replacing both with
validation-based selection is the first item of the revision plan.

\section{Reverse-generation protocol}
\label{app:reverse}

\paragraph{Sampler.} Euler--Maruyama on the reverse SDE of Appendix~\ref{app:notation}, from
$\tmax=3$ down to $\tmin=0.02$, on a geometric time grid (denser at small $t$, where the drift
stiffens like $1/\Dt$). Common random numbers across estimators: the same generator state drives
every arm, so paired differences isolate the score and not the noise.

\paragraph{Why the reference is forward-noised.} Integration stops at $\tmin>0$, so each arm
produces a sample of $P_{\tmin}$, not of $\Pz$. Comparing that against clean data measures the
stopping point. Concretely, adding independent Gaussian noise of variance $v$ to an innovation of
variance $\ivar$ and excess kurtosis $\kappa$ gives excess kurtosis
$\kappa\,(\ivar/(\ivar+v))^2$, since fourth cumulants add and the variance grows by $v$. At
$\corr=0.85$, $\tmin=0.02$ the residual innovation noise is $v = \Dt(1+\corr^2)$ and the
predicted target is $1.9098$ against a clean $3.0$. We therefore draw a reference from $\Pz$ and
noise it forward to the same $\tmin$: both sides are then samples of the same law and the target
follows in closed form. Two earlier versions of this comparison were wrong in exactly the two
available ways --- comparing against clean data, which measured the stopping floor, and applying
a denoising readout, which inverted the bias through shrinkage --- and the closed-form target is
what identified both.

\paragraph{Metrics.} Reported per arm: innovation excess kurtosis with a bootstrap standard error
that resamples chains rather than entries; $\KL(\text{true}\,\|\,\text{generated})$ by histogram;
relative Frobenius error of the covariance and the worst error over lags; marginal variance;
and, on the pointwise side, MSE against the Bayes denoiser with the Bayes risk floor and the
zero-predictor ceiling recorded alongside. Implementation \codefile{src/sample\_metrics.py},
tested by \codefile{tests/test\_sample\_metrics.py}.

\paragraph{Convergence.} Step counts $\{100,200,400,800\}$; see Figure~\ref{fig:steps}. Counts
$1600$ and $3200$ were launched but produced no output and are not reported.

\section{Additional results}
\label{app:extra}

\begin{table}[h]
\caption{Generated innovation excess kurtosis by innovation family, $\nseq=2048$, $\ncomp=4$,
median over three seeds. Families are variance-matched, so they share $\corr^{|i-j|}$ exactly and
differ only beyond second moments. The reference column is grid BP under the true kernel, put
through the same sampler, and is itself a Monte Carlo estimate. The final column is
$|\text{EM--BP}-\text{ref}| - |\text{CNN}-\text{ref}|$: positive means the network is closer.
Source \codefile{outputs/exp\_16/family\_*/generation.csv}. Empirical.}
\label{tab:families}
\centering
\begin{tabular}{lccccr}
\toprule
family & reference & EM--BP & local CNN & global MLP & disagreement \\
\midrule
Gaussian & $-0.022$ & $0.019$ & $0.047$ & $0.053$ & $\mathbf{-0.028}$ \\
mixture  & $-0.467$ & $-0.040$ & $-0.133$ & $0.030$ & $0.094$ \\
uniform  & $-0.766$ & $-0.056$ & $-0.302$ & $0.018$ & $0.245$ \\
Laplace  & $1.921$ & $0.812$ & $1.273$ & $0.170$ & $0.462$ \\
Student  & $2.940$ & $0.989$ & $1.478$ & $0.191$ & $0.489$ \\
\bottomrule
\end{tabular}
\end{table}


\begin{figure}[h]
\centering
\includegraphics[width=0.86\textwidth]{fig_innovation_density.pdf}
\caption{The fitted innovation density against the true Laplace density, at three data budgets
and two capacities, on a log scale so the tails are visible. The deficit at $\ncomp=4$ is in the
tail, not the bulk, which is the mechanism behind the generative disagreement of
\S\ref{sec:generation}. Fitted excess kurtosis is $1.63$--$3.45$ at $\ncomp=4$ and
$2.83$--$3.06$ at $\ncomp=8$ against a true $3.0$; the fitted innovation mean, which is monitored
rather than constrained, falls from $1.9\times10^{-2}$ at $\nseq=128$ to $3.6\times10^{-4}$ at
$\nseq=2048$. Source \codefile{outputs/exp\_18/innovation\_density.csv} and
\codefile{innovation\_summary.csv}. Empirical.}
\label{fig:density}
\end{figure}

\begin{figure}[h]
\centering
\includegraphics[width=0.72\textwidth]{fig_sampler_steps.pdf}
\caption{Reverse-sampler convergence. Generated innovation excess kurtosis and
$\KL(\text{true}\,\|\,\text{generated})$ against the number of Euler--Maruyama steps, all three
arms; error bars are bootstrap standard errors resampling chains. The reference arm sits within
one standard error of the closed-form target at every step count and the arm ordering does not
change, so the comparison of \S\ref{sec:generation} is not an artefact of one discretisation over
this range. Source \codefile{outputs/exp\_16/calibrate\_steps*/steps.csv}. Empirical.}
\label{fig:steps}
\end{figure}

\begin{figure}[h]
\centering
\begin{minipage}[t]{0.48\textwidth}
\centering
\includegraphics[width=\textwidth]{fig_families.pdf}
\end{minipage}\hfill
\begin{minipage}[t]{0.48\textwidth}
\centering
\includegraphics[width=\textwidth]{fig_cost.pdf}
\end{minipage}
\caption{Left: generated-law diagnostics across the five innovation families at $\ncomp=4$ ---
gap to the reference arm's kurtosis, and worst covariance-lag error. The Gaussian family is the
negative control of Table~\ref{tab:families}. Right: measured cost. Inference is $210$--$289$
times slower per sequence than a forward pass of the network; training cost grows linearly in
$\nseq$ for the estimator and is fixed by the step count for the network. Sources
\codefile{outputs/exp\_16/family\_*/generation.csv},
\codefile{outputs/exp\_07\_em\_vs\_score\_network/\{inference,training\}\_cost.csv}.}
\label{fig:extra}
\end{figure}

\begin{figure}[h]
\centering
\includegraphics[width=0.45\textwidth]{fig_closure_families.pdf}
\caption{Error of the Gaussian second-order baseline against grid BP under the true kernel, by
innovation family, at $\corr=0.85$. The ordering by innovation excess kurtosis is visible and the
decay with $t$ is common to all families. Source
\codefile{outputs/exp\_03\_nongaussian\_innovation\_sweep/innovation\_sweep.csv}.}
\label{fig:closure-families}
\end{figure}

\paragraph{Experiments that did not settle their question.} Two are worth recording because they
constrain what may be claimed. The receptive-field sweep (\codefile{exp\_12}) measures how error
falls with the window radius of a local predictor, which is a proxy for the locality of the score
but not a statement about it; without a validation split its selected radius is not defensible,
so its numbers are exploratory. The five-family check that the information-form recursion
reproduces the closed-form Gaussian posterior mean is asserted in
\codefile{tests/test\_gaussian\_bp\_equivalence.py} at $10^{-12}$ for the Laplace chain but is
not persisted as an output for the other families, so no aggregate figure is quoted.

\section{Reproducibility map}
\label{app:repro}

\begin{center}
\scriptsize
\setlength{\tabcolsep}{4pt}
\renewcommand{\arraystretch}{1.15}
\begin{tabular}{@{}>{\raggedright\arraybackslash}p{0.30\textwidth}
                   >{\raggedright\arraybackslash}p{0.30\textwidth}
                   >{\raggedright\arraybackslash}p{0.34\textwidth}@{}}
\toprule
claim & script (\texttt{experiments/\allowbreak }) & output (\texttt{outputs/\allowbreak }) \\
\midrule
grid and domain convergence & \texttt{exp\_\allowbreak 01\_\allowbreak grid\_\allowbreak validation} &
  \texttt{exp\_\allowbreak 01\_\allowbreak .../\allowbreak grid\_\allowbreak heatmap.csv} \\
boundary and quadrature diagnostics & \texttt{exp\_18\_... -{}-parts boundary} &
  \texttt{exp\_\allowbreak 18/\allowbreak boundary.csv} \\
closure error vs $t$ & \texttt{exp\_\allowbreak 02\_\allowbreak laplace\_\allowbreak gaussian\_\allowbreak ...} &
  \texttt{exp\_\allowbreak 02\_\allowbreak .../\allowbreak laplace\_\allowbreak summary.csv} \\
closure error by family & \texttt{exp\_\allowbreak 03\_\allowbreak nongaussian\_\allowbreak ...} &
  \texttt{exp\_\allowbreak 03\_\allowbreak .../\allowbreak innovation\_\allowbreak sweep.csv} \\
EM monotonicity, $\corr$ recovery & \texttt{exp\_18\_... -{}-parts emtrace} &
  \texttt{exp\_\allowbreak 18/\allowbreak em\_\allowbreak trace.csv} \\
learned innovation density & \texttt{exp\_18\_... -{}-parts density} &
  \texttt{exp\_\allowbreak 18/\allowbreak innovation\_\allowbreak density.csv} \\
parameter-recovery rate & \texttt{exp\_\allowbreak 06\_\allowbreak em\_\allowbreak parameter\_\allowbreak ...} &
  \texttt{exp\_\allowbreak 06\_\allowbreak .../\allowbreak em\_\allowbreak rate.csv} \\
Fisher information vs $t$ & \texttt{exp\_\allowbreak 06\_\allowbreak em\_\allowbreak parameter\_\allowbreak ...} &
  \texttt{exp\_\allowbreak 06\_\allowbreak .../\allowbreak price\_\allowbreak of\_\allowbreak noising.csv} \\
Fisher identity vs finite diff. & \texttt{exp\_\allowbreak 08\_\allowbreak gradient\_\allowbreak vs\_\allowbreak ...} &
  \texttt{exp\_\allowbreak 08\_\allowbreak .../\allowbreak gradient\_\allowbreak vs\_\allowbreak exact.csv} \\
sample efficiency (Tab.~\ref{tab:pointwise}) & \texttt{exp\_\allowbreak 07\_\allowbreak em\_\allowbreak vs\_\allowbreak score\_\allowbreak ...} &
  \texttt{replicates/\allowbreak merged\_\allowbreak summary.csv} \\
network capacity control & \texttt{exp\_\allowbreak 07\_\allowbreak em\_\allowbreak vs\_\allowbreak score\_\allowbreak ...} &
  \texttt{exp\_\allowbreak 07\_\allowbreak .../\allowbreak capacity.csv} \\
inference and training cost & \texttt{exp\_\allowbreak 07\_\allowbreak em\_\allowbreak vs\_\allowbreak score\_\allowbreak ...} &
  \texttt{exp\_\allowbreak 07\_\allowbreak .../\allowbreak \{inference,training\}\_cost.csv} \\
receptive field (exploratory) & \texttt{exp\_\allowbreak 12\_\allowbreak receptive\_\allowbreak field} &
  \texttt{exp\_\allowbreak 12\_\allowbreak .../\allowbreak efficiency\_\allowbreak three\_\allowbreak way.csv} \\
sampler step convergence & \texttt{exp\_\allowbreak 16\_\allowbreak ... MODE=calibrate} &
  \texttt{exp\_\allowbreak 16/\allowbreak calibrate\_\allowbreak steps*/\allowbreak steps.csv} \\
generation by budget & \texttt{exp\_\allowbreak 16\_\allowbreak ... MODE=generate} &
  \texttt{exp\_\allowbreak 16/\allowbreak gen\_\allowbreak n*\_\allowbreak seed*/\allowbreak generation.csv} \\
generation by family & \texttt{exp\_\allowbreak 16\_\allowbreak sampling\_\allowbreak validation} &
  \texttt{exp\_\allowbreak 16/\allowbreak family\_\allowbreak */\allowbreak generation.csv} \\
capacity: generation & \texttt{exp\_\allowbreak 16\_\allowbreak sampling\_\allowbreak validation} &
  \texttt{exp\_\allowbreak 16/\allowbreak components\_\allowbreak C*/\allowbreak generation.csv} \\
capacity: pointwise & \texttt{exp\_\allowbreak 16\_\allowbreak sampling\_\allowbreak validation} &
  \texttt{exp\_\allowbreak 16/\allowbreak cpoint\_\allowbreak C*/\allowbreak pointwise.csv} \\
discrete-alphabet variant & \texttt{exp\_\allowbreak 10\_\allowbreak discrete\_\allowbreak alphabet} &
  \texttt{exp\_\allowbreak 10\_\allowbreak .../\allowbreak discrete\_\allowbreak *.csv} \\
\bottomrule
\end{tabular}
\end{center}

Every figure in this note is produced by \codefile{paper/figures/make\_figures.py}, which reads
the outputs above and fails loudly if one is missing; no number in any figure is entered by hand.
Exact commands, environment and HPC job configuration are in \codefile{REPRODUCIBILITY.md}, and
the claim-level audit behind this revision is \codefile{REVISION\_AUDIT.md}.
